% GENERATED FILE. Edit canonical lessons, then run npm run generate:books.
% canonical-source-hash: fnv1a64:332a38f309ab3bc1
% canonical-lessons: TE-C429-purogati, TE-C429-mundadugu, TE-C429-venukabatu, TE-C429-meruguparacu, TE-C429-parinamincu, TE-R429-review-422-429

\chapter{Progress and Setbacks}
\label{ch:te-progress-and-setbacks}
\hlchaptermodality{429}

\begin{chapteropening}
\textbf{By the end of this chapter:} \emph{I can talk about progress, steps forward and improvement.}

Complete the last lesson of chapter 429: I can talk about progress, steps forward and improvement.
\end{chapteropening}

\section[\texorpdfstring{\te{పురోగతి} (purōgati)}{పురోగతి (purōgati)}]{\te{పురోగతి} (purōgati) — progress; headway}

\label{lesson:TE-C429-purogati}



\begin{quote}
Before the new one: say the Telugu for a flaw, then the Telugu for to amend.
\end{quote}

\subsection*{You'll want to know: \te{పురోగతి}}
\textbf{\te{పురోగతి}} — \emph{purōgati} — \textquotedblleft{}progress; headway\textquotedblright{}.

\begin{quote}
\te{ఆ} \te{పనిలో} \te{మంచి} \te{పురోగతి} \te{ఉంది}.
\end{quote}

\begin{quote}
\emph{ā panilō mañci purōgati undi.}
\end{quote}

\begin{quote}
There is good progress in that work.
\end{quote}

Formal; \te{పురోగమించు}, to advance.

\begin{cousinweb}
\textbf{\te{పురోగతి}} (\emph{purōgati}) is built from Sanskrit \emph{puras}, \textquotedblleft{}in front, forward\textquotedblright{}, and \emph{gati}, \textquotedblleft{}going, motion\textquotedblright{} (Monier-Williams; CDIAL 4009).
\end{cousinweb}

\begin{grammarlens}[title={What you've built}]
One more word for saying what happens if.
\end{grammarlens}

\subsection*{Your turn}
\begin{itemize}
\raggedright
  \item \emph{Say it:} \emph{purōgati}
  \item \emph{Say it:} \emph{purōgati}, once more
  \item \emph{Say it:} \emph{ā panilō mañci purōgati undi.} — There is good progress in that work.
  \item \emph{Recall:} say the Telugu for a flaw, then the Telugu for to amend, then say \emph{purōgati} again
\end{itemize}

\begin{tcolorbox}[breakable,colback=teal!4,colframe=teal!35!black,title={Before you move on}]
What does \te{పురోగతి} mean? (\textquotedblleft{}Progress; headway\textquotedblright{}.) Say it once more.
\end{tcolorbox}
\section[\texorpdfstring{\te{ముందడుగు} (mundaḍugu)}{ముందడుగు (mundaḍugu)}]{\te{ముందడుగు} (mundaḍugu) — a step forward}

\label{lesson:TE-C429-mundadugu}



\begin{quote}
Before the new one: say the Telugu for to amend, then the Telugu for progress.
\end{quote}

\subsection*{You'll want to know: \te{ముందడుగు}}
\textbf{\te{ముందడుగు}} — \emph{mundaḍugu} — \textquotedblleft{}a step forward\textquotedblright{}.

\begin{quote}
\te{ఆ} \te{చట్టం} \te{సమానత్వం} \te{వైపు} \te{ఒక} \te{ముందడుగు}.
\end{quote}

\begin{quote}
\emph{ā caṭṭaṁ samānatvaṁ vaipu oka mundaḍugu.}
\end{quote}

\begin{quote}
That law is a step forward towards equality.
\end{quote}

\te{ముందు}, forward, + \te{అడుగు}, a step. \te{ముందడుగు} \te{వేయు}, to take a step forward.

\begin{cousinweb}
\textbf{\te{ముందడుగు}} (\emph{mundaḍugu}) is native Dravidian, a compound of \emph{mundu}, \textquotedblleft{}front, before\textquotedblright{} (DEDR 5020), and \emph{aḍugu}, \textquotedblleft{}a foot, a step\textquotedblright{} (DEDR 72).
\end{cousinweb}

\begin{grammarlens}[title={What you've built}]
One more word for saying what happens if.
\end{grammarlens}

\subsection*{Your turn}
\begin{itemize}
\raggedright
  \item \emph{Say it:} \emph{mundaḍugu}
  \item \emph{Say it:} \emph{mundaḍugu}, once more
  \item \emph{Say it:} \emph{ā caṭṭaṁ samānatvaṁ vaipu oka mundaḍugu.} — That law is a step forward towards equality.
  \item \emph{Recall:} say the Telugu for to amend, then the Telugu for progress, then say \emph{mundaḍugu} again
\end{itemize}

\begin{tcolorbox}[breakable,colback=teal!4,colframe=teal!35!black,title={Before you move on}]
What does \te{ముందడుగు} mean? (\textquotedblleft{}A step forward\textquotedblright{}.) Say it once more.
\end{tcolorbox}
\section[\texorpdfstring{\te{వెనుకబాటు} (venukabāṭu)}{వెనుకబాటు (venukabāṭu)}]{\te{వెనుకబాటు} (venukabāṭu) — backwardness; a setback}

\label{lesson:TE-C429-venukabatu}



\begin{quote}
Before the new one: say the Telugu for progress, then the Telugu for a step forward.
\end{quote}

\subsection*{You'll want to know: \te{వెనుకబాటు}}
\textbf{\te{వెనుకబాటు}} — \emph{venukabāṭu} — \textquotedblleft{}backwardness; a setback\textquotedblright{}.

\begin{quote}
\te{ఆ} \te{గ్రామం} \te{వెనుకబాటుకు} \te{కారణం} \te{విద్య} \te{లేకపోవడం}.
\end{quote}

\begin{quote}
\emph{ā grāmaṁ venukabāṭuku kāraṇaṁ vidya lēkapōvaḍaṁ.}
\end{quote}

\begin{quote}
The reason the village has fallen behind is the lack of education.
\end{quote}

\te{వెనుకబడిన}, backward: \te{వెనుకబడిన} \te{ప్రాంతం}, a backward region.

\begin{cousinweb}
\textbf{\te{వెనుకబాటు}} (\emph{venukabāṭu}) is native Dravidian, built on \emph{venuka}, \textquotedblleft{}the back, the rear\textquotedblright{} (DEDR 5488); the ending -\te{బాటు}, which makes a noun, is of uncertain origin.
\end{cousinweb}

\begin{grammarlens}[title={What you've built}]
One more word for saying what happens if.
\end{grammarlens}

\subsection*{Your turn}
\begin{itemize}
\raggedright
  \item \emph{Say it:} \emph{venukabāṭu}
  \item \emph{Say it:} \emph{venukabāṭu}, once more
  \item \emph{Say it:} \emph{ā grāmaṁ venukabāṭuku kāraṇaṁ vidya lēkapōvaḍaṁ.} — The reason the village has fallen behind is the lack of education.
  \item \emph{Recall:} say the Telugu for progress, then the Telugu for a step forward, then say \emph{venukabāṭu} again
\end{itemize}

\begin{tcolorbox}[breakable,colback=teal!4,colframe=teal!35!black,title={Before you move on}]
What does \te{వెనుకబాటు} mean? (\textquotedblleft{}Backwardness; a setback\textquotedblright{}.) Say it once more.
\end{tcolorbox}
\section[\texorpdfstring{\te{మెరుగుపరచు} (meruguparacu)}{మెరుగుపరచు (meruguparacu)}]{\te{మెరుగుపరచు} (meruguparacu) — to improve; to better}

\label{lesson:TE-C429-meruguparacu}



\begin{quote}
Before the new one: say the Telugu for a step forward, then the Telugu for backwardness.
\end{quote}

\subsection*{You'll want to know: \te{మెరుగుపరచు}}
\textbf{\te{మెరుగుపరచు}} — \emph{meruguparacu} — \textquotedblleft{}to improve; to better\textquotedblright{}.

\begin{quote}
\te{ఆమె} \te{తన} \te{తెలుగును} \te{మెరుగుపరిచింది}.
\end{quote}

\begin{quote}
\emph{āme tana telugunu meruguparicindi.}
\end{quote}

\begin{quote}
She improved her Telugu.
\end{quote}

Past (I …): \te{మెరుగుపరిచాను}, \emph{meruguparicānu}.

\te{మెరుగు}, polish, betterment, + -\te{పరచు}. \te{మెరుగైన}, better.

\begin{cousinweb}
\textbf{\te{మెరుగుపరచు}} (\emph{meruguparacu}) is native Dravidian: \emph{merugu}, \textquotedblleft{}lustre, polish\textquotedblright{} (DEDR 5074, with Kannada \emph{merugu}, \textquotedblleft{}shine\textquotedblright{}, and Telugu \emph{meṟacu}, \textquotedblleft{}to glitter\textquotedblright{}), with -\te{పరచు}, a native ending that makes a verb (DEDR 3853 lists \emph{paṟacu}, \textquotedblleft{}to cause to feel\textquotedblright{}, beside \emph{paḍu}).
\end{cousinweb}

\begin{grammarlens}[title={What you've built}]
One more word for saying what happens if.
\end{grammarlens}

\subsection*{Your turn}
\begin{itemize}
\raggedright
  \item \emph{Say it:} \emph{meruguparacu}
  \item \emph{Say it:} \emph{meruguparacu}, once more
  \item \emph{Say it:} \emph{āme tana telugunu meruguparicindi.} — She improved her Telugu.
  \item \emph{Recall:} say the Telugu for a step forward, then the Telugu for backwardness, then say \emph{meruguparacu} again
\end{itemize}

\begin{tcolorbox}[breakable,colback=teal!4,colframe=teal!35!black,title={Before you move on}]
What does \te{మెరుగుపరచు} mean? (\textquotedblleft{}To improve; to better\textquotedblright{}.) Say it once more.
\end{tcolorbox}
\section[\texorpdfstring{\te{పరిణమించు} (pariṇamiñcu)}{పరిణమించు (pariṇamiñcu)}]{\te{పరిణమించు} (pariṇamiñcu) — to turn into; to develop into}

\label{lesson:TE-C429-parinamincu}



\begin{quote}
Before the new one: say the Telugu for backwardness, then the Telugu for to improve.
\end{quote}

\subsection*{You'll want to know: \te{పరిణమించు}}
\textbf{\te{పరిణమించు}} — \emph{pariṇamiñcu} — \textquotedblleft{}to turn into; to develop into\textquotedblright{}.

\begin{quote}
\te{చిన్న} \te{గొడవ} \te{పెద్ద} \te{సమస్యగా} \te{పరిణమించింది}.
\end{quote}

\begin{quote}
\emph{cinna goḍava pedda samasyagā pariṇamiñcindi.}
\end{quote}

\begin{quote}
A small quarrel turned into a big problem.
\end{quote}

Past (it …): \te{పరిణమించింది}, \emph{pariṇamiñcindi}.

What it becomes takes -\te{గా}; situations do this, so its past is shown as \te{పరిణమించింది}, it turned into. The noun \te{పరిణామం}, a consequence, you know.

\begin{cousinweb}
\textbf{\te{పరిణమించు}} (\emph{pariṇamiñcu}) is built on Sanskrit \emph{pariṇāma}, \textquotedblleft{}change, alteration, transformation\textquotedblright{} (Monier-Williams), with the Telugu verb ending -\te{ఇంచు}.
\end{cousinweb}

\begin{grammarlens}[title={What you've built}]
One more word for saying what happens if.
\end{grammarlens}

\subsection*{Your turn}
\begin{itemize}
\raggedright
  \item \emph{Say it:} \emph{pariṇamiñcu}
  \item \emph{Say it:} \emph{pariṇamiñcu}, once more
  \item \emph{Say it:} \emph{cinna goḍava pedda samasyagā pariṇamiñcindi.} — A small quarrel turned into a big problem.
  \item \emph{Recall:} say the Telugu for backwardness, then the Telugu for to improve, then say \emph{pariṇamiñcu} again
\end{itemize}

\begin{tcolorbox}[breakable,colback=teal!4,colframe=teal!35!black,title={Before you move on}]
What does \te{పరిణమించు} mean? (\textquotedblleft{}To turn into; to develop into\textquotedblright{}.) Say it once more.
\end{tcolorbox}
\section[(dialogue)]{Words from chapters 422–429 — a review}

\label{lesson:TE-R429-review-422-429}



\begin{quote}
Say the Telugu for a guarantee and to fulfil.
\end{quote}

\subsection*{Your turn}
Say these aloud:

\begin{itemize}
\raggedright
  \item a guarantee, to fulfil, to intend, a resolve, to decide firmly
  \item a reward, in vain, needless, to be enough, to make use of
  \item pressure, force, to appeal, to stand up to, to withdraw
  \item a decision, a resolution, to sort out, to finalise, to reconsider
  \item to hold back, to be willing, to turn to, deliberately, a beginning
  \item to insist, to fix up, to reassure, to satisfy, a mediator
  \item to reproach, to insult, to ignore, a flaw, to amend
  \item progress, a step forward, backwardness, to improve, to turn into
\end{itemize}

\begin{tcolorbox}[breakable,colback=teal!4,colframe=teal!35!black,title={Before you move on}]
A guarantee? (\textbf{\te{హామీ}}, \emph{hāmī}.) To fulfil? (\textbf{\te{నెరవేర్చు}}, \emph{neravērcu}.) To intend? (\textbf{\te{ఉద్దేశించు}}, \emph{uddēśiñcu}.) A resolve? (\textbf{\te{సంకల్పం}}, \emph{saṅkalpaṁ}.) To decide firmly? (\textbf{\te{నిశ్చయించు}}, \emph{niścayiñcu}.) A reward? (\textbf{\te{ప్రతిఫలం}}, \emph{pratiphalaṁ}.) In vain? (\textbf{\te{వృథా}}, \emph{vṛthā}.) Needless? (\textbf{\te{అనవసరం}}, \emph{anavasaraṁ}.) To be enough? (\textbf{\te{సరిపోవు}}, \emph{saripōvu}.) To make use of? (\textbf{\te{వినియోగించు}}, \emph{viniyōgiñcu}.) Pressure? (\textbf{\te{ఒత్తిడి}}, \emph{ottiḍi}.) Force? (\textbf{\te{బలవంతం}}, \emph{balavantaṁ}.) To appeal? (\textbf{\te{అభ్యర్థించు}}, \emph{abhyarthiñcu}.) To stand up to? (\textbf{\te{ఎదిరించు}}, \emph{ediriñcu}.) To withdraw? (\textbf{\te{ఉపసంహరించు}}, \emph{upasaṁhariñcu}.) A decision? (\textbf{\te{నిర్ణయం}}, \emph{nirṇayaṁ}.) A resolution? (\textbf{\te{తీర్మానం}}, \emph{tīrmānaṁ}.) To sort out? (\textbf{\te{తేల్చు}}, \emph{tēlcu}.) To finalise? (\textbf{\te{ఖరారు} \te{చేయు}}, \emph{kharāru cēyu}.) To reconsider? (\textbf{\te{పునరాలోచించు}}, \emph{punarālōciñcu}.) To hold back? (\textbf{\te{సంకోచించు}}, \emph{saṅkōciñcu}.) To be willing? (\textbf{\te{సిద్ధపడు}}, \emph{siddhapaḍu}.) To turn to? (\textbf{\te{ఆశ్రయించు}}, \emph{āśrayiñcu}.) Deliberately? (\textbf{\te{ఉద్దేశపూర్వకంగా}}, \emph{uddēśapūrvakaṅgā}.) A beginning? (\textbf{\te{ఆరంభం}}, \emph{ārambhaṁ}.) To insist? (\textbf{\te{పట్టుబట్టు}}, \emph{paṭṭubaṭṭu}.) To fix up? (\textbf{\te{కుదుర్చు}}, \emph{kudurcu}.) To reassure? (\textbf{\te{సమాధానపరచు}}, \emph{samādhānaparacu}.) To satisfy? (\textbf{\te{సంతృప్తిపరచు}}, \emph{santṛptiparacu}.) A mediator? (\textbf{\te{మధ్యవర్తి}}, \emph{madhyavarti}.) To reproach? (\textbf{\te{నిందించు}}, \emph{nindiñcu}.) To insult? (\textbf{\te{అవమానించు}}, \emph{avamāniñcu}.) To ignore? (\textbf{\te{విస్మరించు}}, \emph{vismariñcu}.) A flaw? (\textbf{\te{లోపం}}, \emph{lōpaṁ}.) To amend? (\textbf{\te{సవరించు}}, \emph{savariñcu}.) Progress? (\textbf{\te{పురోగతి}}, \emph{purōgati}.) A step forward? (\textbf{\te{ముందడుగు}}, \emph{mundaḍugu}.) Backwardness? (\textbf{\te{వెనుకబాటు}}, \emph{venukabāṭu}.) To improve? (\textbf{\te{మెరుగుపరచు}}, \emph{meruguparacu}.) To turn into? (\textbf{\te{పరిణమించు}}, \emph{pariṇamiñcu}.)
\end{tcolorbox}
